% =====================================================================
%  landau.tex —— Landau 方程：Feynman 参数表示下的必要条件、动量空间夹点判据、
%  一圈无质量顶角与盒图（第 3 章 / 附录 D 用）
%  研究目录：research/c_regions；所有结论由 DERIVATION.md 所列脚本核对
%  （landau_oneloop.py、vertex_collinear_numeric.py）。
%  依赖：preamble_fact.tex（定理环境、\tagR 等、\ii \dd \sla、费曼图样式）。
% =====================================================================
\section{Landau 方程与夹点}
\label{sec:landau}

\subsection{Feynman 参数表示}

考虑 $L$ 圈、$N$ 条内线的图 $G$。第 $j$ 条线的动量 $q_j=\sum_l\eta_{jl}k_l+r_j$（$\eta_{jl}\in\{0,\pm1\}$，$r_j$ 为外动量
的线性组合），质量 $m_j$，传播子 $\ii/(D_j+\ii0)$，$D_j\equiv q_j^2-m_j^2$。对标量被积函数，
\begin{multline}\label{eq:landau-param}
\int\prod_{l=1}^{L}\frac{\dd^Dk_l}{\ii\pi^{D/2}}\prod_{j=1}^{N}\frac{1}{(-D_j-\ii0)^{\nu_j}}\\
=\frac{\Gamma(\nu-LD/2)}{\prod_j\Gamma(\nu_j)}\int_{\mathbb R_+^N}\prod_j\dd\alpha_j\,\alpha_j^{\nu_j-1}\,
\delta\Big(1-\sum_{j\in\sigma}\alpha_j\Big)\frac{\mathcal U^{\,\nu-(L+1)D/2}}{(\mathcal F-\ii0)^{\nu-LD/2}},
\end{multline}
$\nu=\sum_j\nu_j$，$\sigma$ 为任意非空子集（Cheng--Wu 定理）。记二次型
$\mathcal D(\alpha,k)\equiv\sum_j\alpha_j(m_j^2-q_j^2)$，其对 $k$ 的驻点为 $k_c(\alpha)$；则
$\mathcal U=\det A=\sum_{T\in\mathcal T_1}\prod_{j\notin T}\alpha_j$，
$\mathcal F=\mathcal U\,\mathcal D(\alpha,k_c)=-\sum_{T\in\mathcal T_2}\big(\prod_{j\notin T}\alpha_j\big)s_T+\mathcal U\sum_j\alpha_jm_j^2$，
其中 $\mathcal T_{1}$、$\mathcal T_{2}$ 为生成树与 2-森林，$s_T$ 为流入 2-森林一个连通分支的动量的平方。
$\mathcal U$、$\mathcal F$ 分别是 $L$、$L+1$ 次齐次多项式。

\subsection{参数空间形式的 Landau 条件}

\begin{theorem}[Landau 条件，参数形式]\label{thm:landau-alpha}\tagR\
设 $p_0$ 为实外运动学，$D$ 取值使 \cref{eq:landau-param} 中除 $\mathcal F$ 的零点外再无不可积性（紫外收敛假设）。
若闭单形 $\bar\Delta=\{\alpha_j\ge0,\ \sum_j\alpha_j=1\}$ 上不存在满足
\begin{equation}\label{eq:landau-alpha}
\mathcal F(\alpha,p_0)=0,\qquad \alpha_j\,\frac{\partial\mathcal F}{\partial\alpha_j}(\alpha,p_0)=0\quad(j=1,\dots,N)
\end{equation}
的点，则积分在 $p_0$ 处有限，且作为不变量的函数可全纯延拓到 $p_0$ 的复邻域。等价地：$p_0$ 处的红外发散或奇点
必伴随 \cref{eq:landau-alpha} 的解。由 Euler 关系 $(L+1)\mathcal F=\sum_j\alpha_j\partial_j\mathcal F$，第一式由其余各式推出。
\end{theorem}

\begin{proof}
取 $\sigma=\{1,\dots,N\}$。令 $c(\alpha)=\sum_k\alpha_k\partial_k\mathcal F$，$w_j=\partial_j\mathcal F-c$，作形变
$z_j(\alpha,\tau)=\alpha_j\,[1-\ii\tau\eta\,w_j(\alpha)]$，$0\le\tau\le1$，$\eta>0$ 充分小。
(i) $\sum_jz_j=1-\ii\tau\eta(c-c)=1$，约束保持。
(ii) $\alpha_j=0$ 时 $z_j=0$：单形的各面被映入相应的复面，Stokes 公式中“侧面”的贡献为零
（全纯 $(N-1,0)$ 形式限制到落在 $(N-2)$ 维复面内的实链上为零）；因而参数空间的端点奇异与内部奇异被同样处理。
(iii) 对实 $p$，$\mathcal F(z)=\mathcal F(\alpha)-\ii\tau\eta V(\alpha)+O(\eta^2)$，其中
$V=\sum_j\alpha_jw_j^2=\sum_j\alpha_j(\partial_j\mathcal F)^2-c^2\ge0$（用到 $\sum_j\alpha_jw_j=0$），且 $O(\eta^2)$ 项为实数，故
$\Imag\mathcal F(z)=-\tau\eta V+O(\eta^3)$。
(iv) 若 $V=\mathcal F=0$，则凡 $\alpha_j>0$ 处 $\partial_j\mathcal F=c$，于是 $\mathcal F=c/(L+1)=0$，即 $c=0$、$\alpha_j\partial_j\mathcal F=0$：
这正是 \cref{eq:landau-alpha} 的解。按假设无解，故连续函数 $|\mathcal F|+V$ 在 $\bar\Delta\times\mathcal N$ 上有正下界 $2\delta$
（$\mathcal N$ 为 $p_0$ 的邻域，紧性）。$\eta$ 足够小时：$|\mathcal F(\alpha)|\ge\delta$ 处 $|\mathcal F(z)|>0$；$|\mathcal F(\alpha)|<\delta$ 处
$V\ge\delta$，$\Imag\mathcal F(z)<0$。所以对一切 $\tau\in[0,1]$、$\varepsilon\ge0$ 均有 $\mathcal F(z)-\ii\varepsilon\ne0$：同伦从不在 $-\ii0$
所禁止的一侧穿过奇异集，原积分的 $\varepsilon\to0^+$ 极限等于形变后的积分。
(v) $\mathcal U(z)$ 是正单项式乘以落在正实轴附近窄锥内的因子之和，$|\mathcal U(z)|\ge\cos\theta\,\mathcal U(\alpha)$，紫外可积性不受影响。
(vi) 形变后的被积函数有界于 $C\,\mathcal U(\alpha)^{a}\prod\alpha_j^{\nu_j-1}$ 且对复化的 $p\in\mathcal N$ 全纯，由控制收敛定理即得结论。
\end{proof}

\begin{lemma}[两种形式的等价]\label{lem:landau-equiv}\tagR\
设 $\mathcal U(\alpha)\neq0$。则 \cref{eq:landau-alpha} 成立当且仅当存在实动量 $k$ 使
\begin{equation}\label{eq:landau-momentum}
\alpha_j\,D_j(k)=0\quad(\forall j),\qquad \sum_j\alpha_j\,\eta_{jl}\,q_j^\mu(k)=0\quad(\forall l,\ \text{圈方程}),
\end{equation}
且此时 $k=k_c(\alpha)$。
\end{lemma}
\begin{proof}
圈方程即 $\partial\mathcal D/\partial k_l=0$，故 $k=k_c$。包络定理给出
$\partial_j(\mathcal F/\mathcal U)=\partial\mathcal D/\partial\alpha_j|_{k_c}=-D_j(k_c)$。在 $\mathcal F=0$ 处 $\partial_j\mathcal F=\mathcal U\,\partial_j(\mathcal F/\mathcal U)$，
故 $\alpha_j\partial_j\mathcal F=-\mathcal U\alpha_jD_j(k_c)$。反之 $\alpha_jD_j(k_c)=0$ 给出 $\mathcal F/\mathcal U=-\sum_j\alpha_jD_j(k_c)=0$。
\end{proof}

\begin{remark}
\cref{lem:landau-equiv} 只是两组\emph{方程}的代数等价，并不意味着“动量空间积分围道被夹住”当且仅当“参数积分围道被夹住”。
Collins（arXiv:2007.04085，附录 D1）指出这一被 Landau 与 Coleman--Norton 默认的蕴涵对无质量图不成立：一圈无质量自能在
$p^2=0$ 时的共线夹点在参数积分中没有对应的夹点（此时 $\mathcal F=-p^2\alpha_1\alpha_2$ 恒为零）。因此
\cref{thm:landau-alpha} 用于外参数的奇异性与红外有限性判据；因子化所需的“圈动量空间中被夹住的区域”由下面的
\cref{thm:landau-pinch} 刻画。
\end{remark}

\subsection{动量空间的夹点判据}

\begin{theorem}[被夹点的必要条件]\label{thm:landau-pinch}\tagR\
固定实外动量。对实点 $k_0\in\mathbb R^{4L}$，令 $S=\{j:D_j(k_0)=0\}$，$u_j=\partial D_j/\partial k\,(k_0)$。若不存在
$\alpha\in\mathbb R^{S}_{\ge0}\setminus\{0\}$ 使 $\sum_{j\in S}\alpha_ju_j=0$，则存在实方向 $v$ 与 $k_0$ 的邻域，使形变
$k\to k+\ii\eta\chi(k)v$（$\chi$ 为光滑截断函数）在 $\chi>0$ 处对 $j\in S$ 给出 $\Imag D_j>0$，并使 $j\notin S$ 的 $D_j$ 在形变后围道上
不为零：$k_0$ 未被夹住。等价地，被夹住的实点满足 $\alpha_j\ge0$ 的动量空间 Landau 方程 \cref{eq:landau-momentum}。
\end{theorem}
\begin{proof}
Gordan 择一定理：对向量组 $\{u_j\}$，下列二者恰有其一成立——(I) 存在 $v$ 使所有 $u_j\cdot v>0$；
(II) 存在 $\alpha\ge0$、$\alpha\ne0$ 使 $\sum\alpha_ju_j=0$。〔若 $0\in\operatorname{conv}\{u_j\}$ 则 (II)；否则取紧凸集
$\operatorname{conv}\{u_j\}$ 中距原点最近的点 $v$，则其上 $u\cdot v\ge|v|^2>0$，即 (I)；二者并存将给出 $0=\sum\alpha_ju_j\cdot v>0$。〕
由假设 (I) 成立。因 $D_j$ 为二次式，$D_j(k+\ii\eta\chi v)=D_j(k)-\eta^2\chi^2(v'_j)^2+2\ii\eta\chi\,q_j(k)\cdot v'_j$ 严格成立
（$v'_j=\sum_l\eta_{jl}v_l$），故 $\Imag D_j=\eta\chi\,u_j(k)\cdot v$，由连续性在 $k_0$ 附近为正；$j\notin S$ 时 $|D_j|$ 在小邻域内有正下界，
而形变只改变 $O(\eta)$。
\end{proof}

\begin{remark}
(a) 零动量无质量线 $u_j=2q_j=0$，总是平凡地满足 Landau 方程；它是否导致发散由幂次计数决定（\cref{sec:powercount}）。
(b) 反方向（Landau 方程 $\Rightarrow$ 被夹住）以及由局部形变拼成整体形变，需排除一阶位移为零的“反常形变”，
完整证明见 Collins 2020（arXiv:2007.04085，INSPIRE 上无期刊发表记录），其适用范围包括标准 Feynman 图及含 Wilson 线
（线性分母）的图；在本书中记为 \tagA。
(c) $\alpha_j\ge0$ 的符号与 $+\ii0$ 规定相联系。
\end{remark}

\paragraph{Coleman--Norton 图像。}
对 $\alpha_j\ge0$ 的实解，令每条线的时空位移 $\Delta x_j=\alpha_jq_j$。圈方程恰是“每个圈上 $\sum\Delta x_j=0$”，即存在顶点位置
使 $x_{\text{终}}-x_{\text{始}}=\alpha_jq_j$：$\alpha_j>0$ 的在壳线是以速度 $\vec q_j/q_j^0$ 经典传播的粒子（无质量线沿光线），$\alpha_j=0$ 的线长度
为零（收缩为有效顶点）。零动量线是退化的：在 $\alpha_j\to\infty$、$\alpha_jq_j$ 有限的标度极限下，它连接任意两点（软交换）。
这一代数改写是平凡的 \tagR；而“被夹住 $\Leftrightarrow$ 可经典实现”还需要 \cref{thm:landau-pinch} 的逆命题（Coleman \& Norton 1965；
其证明中的缺口见 Collins 2020）。

\subsection{例：一圈无质量顶角}

\begin{figure}[htbp]
\centering
\begin{subfigure}[b]{0.31\textwidth}\centering
\begin{tikzpicture}[fd,scale=1.0]
  \coordinate (q) at (-1.7,0); \coordinate (X) at (0,0);
  \coordinate (Y1) at (1.5,0.95); \coordinate (Y2) at (1.5,-0.95);
  \coordinate (o1) at (2.6,1.6); \coordinate (o2) at (2.6,-1.6);
  \draw[photon] (q) -- (X);
  \draw[quark] (X) -- (Y1); \draw[quark] (Y1) -- (o1);
  \draw[quark] (o2) -- (Y2); \draw[quark] (Y2) -- (X);
  \draw[gluon] (Y2) -- (Y1);
  \foreach \v in {X,Y1,Y2} \node[vertex] at (\v) {};
  \node[above] at (-0.9,0.05) {$q$};
  \node[above left] at (0.75,0.45) {$p_1-k$};
  \node[below left] at (0.75,-0.45) {$p_2+k$};
  \node[right] at (1.62,0) {$k$};
  \node[right] at (o1) {$p_1$}; \node[right] at (o2) {$p_2$};
  \node[below=2pt] at (X) {$X$}; \node[above] at (Y1) {$Y_1$}; \node[below] at (Y2) {$Y_2$};
\end{tikzpicture}
\caption{顶角图与动量记号}\label{fig:landau-vertex}
\end{subfigure}\hfill
\begin{subfigure}[b]{0.32\textwidth}\centering
\begin{tikzpicture}[scale=2.3,line cap=round]
  \coordinate (A0) at (0,0); \coordinate (A1) at (1,0); \coordinate (A2) at (0.5,0.866);
  \fill[gray!8] (A0) -- (A1) -- (A2) -- cycle;
  \draw[line width=2.2pt,red!70!black] (A0) -- (A1);
  \draw[line width=2.2pt,blue!70!black] (A0) -- (A2);
  \draw[dashed,gray!70!black,line width=0.8pt] (A1) -- (A2);
  \fill (A0) circle (0.035); \fill[red!70!black] (A1) circle (0.03); \fill[blue!70!black] (A2) circle (0.03);
  \node[below left] at (A0) {\small $S$};
  \node[below right] at (A1) {\small $\alpha_1{=}1$};
  \node[above] at (A2) {\small $\alpha_2{=}1$};
  \node[below,red!70!black] at (0.5,0) {\small $C_1$: $\alpha_2=0$};
  \node[left,blue!70!black] at (0.2,0.5) {\small $C_2$};
  \node[right,gray!50!black,align=left] at (0.78,0.5) {\small $\alpha_0=0$\\[-2pt]\small ($Q^2=0$ 时)};
  \node at (0.45,0.3) {\small $H$};
\end{tikzpicture}
\caption{参数单形上的 Landau 集}\label{fig:landau-simplex}
\end{subfigure}\hfill
\begin{subfigure}[b]{0.33\textwidth}\centering
\begin{tikzpicture}[fd,scale=0.9]
  % hard point
  \begin{scope}[shift={(0,0)}]
    \draw[plain] (0,0) -- (1.2,0.8); \draw[plain] (0,0) -- (1.2,-0.8);
    \node[vertex] at (0,0) {}; \node[below] at (0.55,-0.95) {\small $H$};
  \end{scope}
  % collinear
  \begin{scope}[shift={(1.75,0)}]
    \draw[plain] (0,0) -- (1.2,-0.8);
    \draw[plain] (0,0) -- (1.4,0.93);
    \draw[gluon] (0.02,0.10) -- (0.83,0.64);
    \node[vertex] at (0,0) {}; \node[vertex] at (0.85,0.567) {};
    \node[below] at (0.6,-0.95) {\small $C_1$};
  \end{scope}
  % soft
  \begin{scope}[shift={(3.5,0)}]
    \draw[plain] (0,0) -- (1.3,0.87); \draw[plain] (0,0) -- (1.3,-0.87);
    \draw[gluon] (0.9,0.6) -- (0.7,-0.47);
    \node[vertex] at (0,0) {}; \node[below] at (0.6,-0.95) {\small $S$};
  \end{scope}
\end{tikzpicture}
\caption{Coleman--Norton 时空图像}\label{fig:landau-cn}
\end{subfigure}
\caption{一圈无质量顶角 $\gamma^*\to q\bar q$ 的夹点结构。(b) 中红、蓝两条闭棱为参数空间 Landau 解：$\alpha_2=0$
对应 $k=xp_1$（共线于 $p_1$），$\alpha_1=0$ 对应 $k=-yp_2$；二者交于软点 $S$（$k=0$）；第三条棱 $\alpha_0=0$（$q\bar q$ 两粒子截口）
仅在 $Q^2=0$ 时是解；内部为硬区域 $H$。}
\label{fig:landau-vertex-all}
\end{figure}

取 $\gamma^*(q)\to q(p_1)\bar q(p_2)$，$p_1^2=p_2^2=0$，$q^2=2p_1\cdot p_2=Q^2>0$；内线 $\ell_0=k$（胶子），$\ell_1=p_1-k$，$\ell_2=p_2+k$
（\cref{fig:landau-vertex}）。则
\begin{equation}\label{eq:landau-vertexF}
\mathcal U=\alpha_0+\alpha_1+\alpha_2,\qquad \mathcal F=-Q^2\alpha_1\alpha_2,\qquad k_c=\alpha_1p_1-\alpha_2p_2\quad\Big(\sum\alpha=1\Big).
\end{equation}
\cref{eq:landau-alpha} 给出 $\alpha_1\alpha_2=0$：解集是单形的两条闭棱（\cref{fig:landau-simplex}）。由 \cref{lem:landau-equiv}：
\begin{itemize}[leftmargin=2em]
\item 棱 $\alpha_2=0$（$\alpha_1=x\in[0,1]$）：$k=xp_1$，线 0、1 在壳且共线于 $p_1$，反夸克线离壳——\textbf{共线夹点面} $C_1$；
  圈方程 $\alpha_2(p_2+k)+\alpha_0k+\alpha_1(k-p_1)=0$ 给出 $\alpha_0x=\alpha_1(1-x)$，故物理区 $0\le x\le1$。
\item 棱 $\alpha_1=0$：$k=-yp_2$，$0\le y\le1$——共线夹点面 $C_2$。
\item 顶点 $(1,0,0)$：$k=0$，三条线均在壳但只有 $\alpha_0\ne0$——\textbf{软点} $S$。顶点 $(0,1,0)$、$(0,0,1)$：零动量夸克线，
  即 $C_1$、$C_2$ 的“软夸克”端点（幂次压低）。
\item 棱 $\alpha_0=0$ 与内部：$Q^2\ne0$ 时无解——\textbf{硬点} $H$（所有线收缩为一点）。类时 $Q^2$ 下两夸克线同时在壳的
  截口曲面存在，但按 \cref{thm:landau-pinch} 未被夹住，只产生吸收部分；仅在 $Q^2=0$（法向阈）时被夹住。
\end{itemize}
同样的结论也可由一圈 Cayley 矩阵 $Y_{ij}=q_i\cdot q_j|_{\text{在壳}}$ 直接得到：$Y=\begin{pmatrix}0&0&0\\0&0&-Q^2/2\\0&-Q^2/2&0\end{pmatrix}$，
$Y_S$ 有非负零向量当且仅当 $S\in\{\{0\},\{1\},\{2\},\{0,1\},\{0,2\}\}$。时空图像（\cref{fig:landau-cn}）：$C_1$ 中光子在 $X$ 处产生
反夸克与一对共线的夸克–胶子，二者沿 $p_1$ 的光线同行并在 $Y_1=X+\alpha_1(1-x)p_1=X+\alpha_0xp_1$ 处同时到达而复合；$S$ 中零动量
胶子连接两条出射光线上的任意两点。$C_1$、$C_2$、$S$ 均为对数发散（\cref{sec:powercount}），在 $S$ 处重叠，产生 $1/\epsilon^2$。

\paragraph{无质量盒图。} 取 $q_j=k+r_j$，$r_0=0$，$r_1=p_1$，$r_2=p_1+p_2$，$r_3=-p_4$，$p_i^2=0$，
$\mathcal F=-s\,\alpha_0\alpha_2-t\,\alpha_1\alpha_3$。Landau 集为 3-单形的四条棱 $(01),(12),(23),(30)$（相邻两线在壳且共线于其间流入的外动量）
及四个顶点（零动量线：相邻两外线间的软交换）；棱 $(02)$、$(13)$（$s$、$t$ 道两粒子截口）仅在 $s=0$、$t=0$ 时为解。
标量盒图在 $D=4$ 中每个顶点与每条棱的幂次计数均为对数型，由此给出熟知的 $1/\epsilon^2$ 红外极点。

\begin{remark}[必要而非充分]
无质量图总有平凡的顶点解。对离壳（欧氏）顶角 $\mathcal F=P^2\alpha_0(\alpha_1+\alpha_2)+Q^2\alpha_1\alpha_2$，Landau 集只含三个顶点；
在软顶点附近 $\mathcal F=rP^2+O(r^2)$（$r=\alpha_1+\alpha_2$），$\int r\,\dd r/r$ 收敛，不发散。对有质量在壳外线、无质量胶子，
$\mathcal F=Q^2\alpha_1\alpha_2+m^2(\alpha_1+\alpha_2)^2$，同一顶点仍是唯一解，但 $\mathcal F=r^2(m^2+u(1-u)Q^2)$，$\int r\,\dd r/r^2$ 对数发散
（软红外发散）。两种情形的 Landau 解完全相同，是幂次计数作出区分。
\end{remark}
